\section*{Capítulo \ref{ch:g4:volcanoes-earthquakes-first-look} --- Vulcões e terremotos: um primeiro olhar}

\begin{solution}{exo:g4:volcanoes-earthquakes-first-look:1}
(a) \omterm{def:g4:volcanoes-earthquakes-first-look:eruption}{Lava}. (b) \omterm{def:g4:volcanoes-earthquakes-first-look:eruption}{Cinza vulcânica}. (c) Gás.
\end{solution}

\begin{solution}{exo:g4:volcanoes-earthquakes-first-look:2}
Um \omterm{def:g4:volcanoes-earthquakes-first-look:earthquake}{terremoto} é um tremor repentino do chão, causado por \omterm{def:g3:stones-pebbles-sand:rock}{rochas} que se
quebram e escorregam lá no fundo da terra. Ele dura segundos (um ou dois
minutos para os maiores de todos), e não horas.
\end{solution}

\begin{solution}{exo:g4:volcanoes-earthquakes-first-look:3}
Abaixe-se, proteja-se, segure-se.
\end{solution}

\begin{solution}{exo:g4:volcanoes-earthquakes-first-look:4}
Em volta do \omterm{def:g2:water-on-earth:ocean}{oceano} Pacífico.
\end{solution}

\begin{solution}{exo:g4:volcanoes-earthquakes-first-look:5}
(a) Verdadeiro: só os instrumentos percebem a maioria deles. (b) Falso: a
\omterm{def:g4:volcanoes-earthquakes-first-look:eruption}{lava} esfria e endurece, formando \omterm{def:g3:stones-pebbles-sand:rock}{rocha} nova. (c) Falso: correr enquanto o
chão treme é perigoso (você pode cair, e vidros e objetos caem dos
prédios); abaixe-se, proteja-se e segure-se.
\end{solution}

\begin{solution}{exo:g4:volcanoes-earthquakes-first-look:6}
Ela vem das profundezas, embaixo do \omterm{def:g4:volcanoes-earthquakes-first-look:volcano}{vulcão}, onde a \omterm{def:g3:stones-pebbles-sand:rock}{rocha} derreteu. Ela
sobe pelo meio do \omterm{def:g4:volcanoes-earthquakes-first-look:volcano}{vulcão} e sai no alto.
\end{solution}

\begin{solution}{exo:g4:volcanoes-earthquakes-first-look:7}
(a) O Chile (o Peru, o Equador ou o México também estão certos): os seus
\omterm{def:g4:volcanoes-earthquakes-first-look:volcano}{vulcões} se estendem por todo o seu comprimento, perto da costa do
Pacífico. (b) O Havaí, no meio do \omterm{def:g2:water-on-earth:ocean}{oceano} Pacífico, longe do anel sombreado.
\end{solution}

\begin{solution}{exo:g4:volcanoes-earthquakes-first-look:8}
$141 \div 12$ é um pouco menos que 12 (porque $12 \times 12 = 144$): cerca
de 12 \omterm{def:g4:volcanoes-earthquakes-first-look:earthquake}{terremotos} grandes por mês.
\end{solution}

\begin{solution}{exo:g4:volcanoes-earthquakes-first-look:9}
De 20, um quarto é 5; então três quartos são $3 \times 5 = 15$: cerca de 15
no anel, e cerca de $20 - 15 = 5$ fora dele.
\end{solution}

\begin{solution}{exo:g4:volcanoes-earthquakes-first-look:10}
Ele deveria ter se abaixado, ficado de quatro, protegido a cabeça e o
pescoço embaixo da carteira e se segurado nela até o tremor parar. O vão
da porta não o protege dos objetos que caem ou voam, e correr enquanto o
chão treme pode fazê-lo cair.
\end{solution}

\begin{solution}{exo:g4:volcanoes-earthquakes-first-look:11}
Fora do anel: $1214 - 837 = 377$ \omterm{def:g4:volcanoes-earthquakes-first-look:volcano}{vulcões}. Um terço de 1\,200 é 400; então
dois terços de 1\,200 são 800. O anel tem 837, um pouco mais que 800 em
cerca de 1\,200: ``cerca de 2 em 3'' é uma boa descrição.
\end{solution}
